\clearpage\NoSourceHead\SourceLeaf{414}
\SourceGroup{VI}
\HistoricTitle{BESTIMMUNG DER CONVERGENZ DER REIHEN, IN WELCHE DIE PERIODISCHEN FUNCTIONEN EINER VERÄNDERLICHEN GRÖSSE ENTWICKELT WERDEN.}
\HistoricCitation{[Handschrift von {\scriptsize 16} mit B{\scriptsize 1} bis B{\scriptsize 16} bezeichneten Seiten in Fa, Kapsel {\scriptsize 46}a.]}
\begin{center}1.\end{center}
\noindent\hspace*{\parindent}Die Überschrift bezeichnet zwar den Hauptgegenstand dieser Denkschrift, der jedoch nur einen Theil ihres Inhalts ausmacht. Die Lösung der Aufgabe gehört recht eigentlich in das Gebiet der Lehre von den complexen (imaginären) Grössen, auf welche man in frühern Zeiten die Geschäfte der Analysis nur ausnahmsweise auszudehnen sich erlaubte. Die vollkommene Befugniss, der Analysis gleichmässig das ganze Gebiet der complexen Grössen zu unterwerfen, habe ich im Jahre 1831 nachgewiesen: allein die Wege in diesem Gebiete sind noch nicht überall gebahnt, und die in der Analysis bisher gangbaren Begriffsabgrenzungen sind ursprünglich fast immer nur unter der stillschweigenden Voraussetzung gemacht, dass man über das Gebiet der reellen Grössen nicht hinausgehe, und ermangeln selbst unter dieser Beschränkung nicht selten einer befriedigenden Schärfe. Es werden daher mit dem Hauptgegenstande einige Nebenuntersuchungen und Erörterungen verbunden werden müssen, ohne jedoch bei denselben eine grössere Ausführlichkeit zu beabsichtigen, als für unsere Zwecke nöthig ist.
\begin{center}2.\end{center}
\noindent\hspace*{\parindent}Zuvörderst ist die bekannte Versinnlichung der complexen Grössen in Erinnerung zu bringen, wo jede complexe Grösse \(x+iy\) (die imaginäre Einheit
\clearpage\SourceLeaf{415}\ArithmeticHead{408}{L}{ANALYSIS. NACHLASS.}
\noindent \(\surd{-1}\) durch \(i\) bezeichnet) durch denjenigen Punkt einer unbegrenzten Ebene repräsentirt wird, dessen rechtwinkelige Coordinaten in Beziehung auf ein beliebig gewähltes System, \(x\) und \(y\) sind. Es wird damit, wie schon an einem anderen Orte[\textsuperscript{*})] bemerkt ist, nur bezweckt, die Bewegung in dem an sich vom Räumlichen unabhängigen Felde der abstracten complexen Grössen zu erleichtern und eine Sprache für dasselbe zu vermitteln.

Auf diese Weise versinnlicht sich durch eine Linie in der Ebene eine durch unendlich kleine Differenzen nach der Stetigkeit fortschreitende Reihe complexer Grössen, die ich einen \textls[120]{Zug} nennen werde, so wie jede der einem bestimmten Zuge angehörenden Grössen eine \textls[120]{Stelle} in diesem Zuge. Ist die letzte der einen Zug bildenden complexen Grössen der ersten gleich, während alle dazwischen liegenden unter sich ungleich sind, so wird dieser \textls[120]{geschlossene} Zug vertreten durch eine geschlossene Linie, und der von dieser Linie eingeschlossene Flächenraum vertritt wieder alle diejenigen nach der Stetigkeit mit einander zusammenhängenden complexen Grössen, die in jenem Zuge ihre Limite finden, und die ich eine \textls[120]{Schicht} complexer Grössen nennen werde. Dieser Begriff lässt sich noch allgemeiner auffassen als der Inbegriff aller nach der Stetigkeit unter sich zusammenhängender complexer Grössen, die von den übrigen nicht dazu zu rechnenden durch \textls[120]{einen oder mehrere} geschlossene Züge geschieden sind. In letzterem Falle bildet einer dieser Züge, der äusserste, die Scheidung von dem übrigen unendlichen Gebiete (welches dann auch als eine unendliche Schicht bezeichnet werden mag), die anderen dienen dazu, im Innern eine oder mehrere Schichten --- in der räumlichen Vorstellungsart: Inseln oder Enclaven --- abzusondern.
\begin{center}3.\end{center}
\noindent\hspace*{\parindent}Zu einem Zuge complexer Grössen steht eine dadurch begrenzte Schicht in einer gewissen Beziehung analog derjenigen, wonach von einem Flächenstück ausgesagt wird, dass es der Begrenzungslinie zur Rechten oder zur Linken liege\textsuperscript{**}). Bestimmt ist diese Unterscheidung nur, insofern man in der
\par\vfill\noindent\rule{.25\linewidth}{.4pt}\par\smallskip
{\footnotesize
\noindent\hspace*{\parindent}[\textsuperscript{*}) Siehe den art. {\scriptsize 5}. der Jubiläumsschrift von {\scriptsize 1849} \textit{Beiträge zur Theorie der algebraischen Gleichungen}, Werke III, S. {\scriptsize 79}.]\par
\noindent\hspace*{\parindent}\textsuperscript{**}) Da hier nicht der Ort ist in die Metaphysik des Räumlichen weiter einzudringen, so übergehe ich den Umstand, dass die Wechselbeziehung zwischen vorwärts-\par}
\clearpage\SourceLeaf{416}\ArithmeticHead{409}{R}{CONVERGENZ DER ENTWICKLUNG PERIODISCHER FUNCTIONEN.}
\noindent Linie eine bestimmte Richtung als eine vorwärts gehende annimmt: wählt man dafür die entgegengesetzte, so wird, was vorher rechts war, zu links und umgekehrt. Auf gleiche Weise muss in einem eine Schicht complexer Grössen \(S\) begrenzenden, oder, was dasselbe ist, zwei solcher Schichten \(S\) und \(S'\) scheidenden Grössenzuge eine bestimmte Folgeordnung angenommen werden und das Charakteristische im Verhältniss einer Schicht zu dem Grenzzuge kann dann am einfachsten auf folgende Art ausgesprochen werden: Es seien \(t,t+\alpha\) zwei unendlich wenig verschiedene complexe Grössen in dem Zuge, die erste vorangehend, die andere folgend; \(t+\beta\) eine andere gleichfalls unendlich wenig von \(t\) verschiede[ne] complexe Grösse und
\[
\frac\beta\alpha=k+il.
\]
\noindent Es wird dann, (wenn wir den singulären Fall ausschliessen, wo die Grösse \(t\) in räumlicher Darstellung des Zuges eine Spitze darbietet) \(l=0\) oder unendlich klein sein, wenn \(t+\beta\) dem Zuge selbst angehört; hingegen wird für einen positiven endlichen Werth von \(l\) die Grösse \(t+\beta\) einer der Schichten \(S\) und \(S'\), und für einen negativen der andern angehören; und dies Verhältniss wird für den ganzen Zug dasselbe bleiben, eben so wie der, welcher einen See hart am Ufer so zu umschreiten anfängt, dass jener ihm zur Rechten liegt, denselben fortwährend zur Rechten behält, wenn er nur sich selbst nicht umwendet.

Man kann dies dadurch ausdrücken, dass man sagt, diejenige Schicht, welcher die positiven Werthe von \(l\) entsprechen, habe ein positives Verhältniss zu dem Grenzzuge, die andere ein negatives. Offenbar vertauschen sich diese Benennungen, sobald man in dem Grenzzuge die entgegengesetzte Folgeordnung wählt. Man erkennt hieraus die Bedeutung eines Ausdrucks, dessen ich mich in der Folge zuweilen bedienen werde. Es seien \(F,F'\) die beiden denkbaren Folgeordnungen in einem Zuge \(Z\), der die Schicht \(S\) begrenzt; ebenso
\par\vfill\noindent\rule{.25\linewidth}{.4pt}\par\smallskip
{\footnotesize\noindent rückwärts und rechts-links erst durch die Hinzufügung des dritten Gegensatzes oben-unten zwischen den Raumtheilen, welche die Fläche scheidet, Haltung bekommt, so wie den, dass dieses Verhältniss nicht durch eine Definition a priori gegeben, sondern nur durch Zusammenhalten mit einem \textls[120]{wirklich Vorhandenen}, drei Dimensionen darbietenden erkannt werden kann, insofern in diesem die Namen bereits feststehen. Man vergl. Göttingische gelehrte Anzeigen 1831, S. 637 [Werke II, S. 177]. Die abstracte allgemeine Lehre von den complexen Grössen hat mit beiden nichts zu schaffen.\par}
\smallskip\noindent{\small X1.}\hfill{\small 52}\par
\clearpage\SourceLeaf{417}\ArithmeticHead{410}{L}{ANALYSIS. NACHLASS.}
\noindent \(f,f'\) die beiden Folgeordnungen in einem Zuge \(z\), welcher die Schicht \(s\) begrenzt. Ist nun das Verhältniss von \(S\) zu \(Z\), unter Voraussetzung der Folgeordnung \(F\) dasselbe wie das Verhältniss von \(s\) zu \(z\) unter Voraussetzung der Folgeordnung \(f\), d. i. sind beide positiv oder beide negativ, so drücke ich dies dadurch aus, dass ich die Folgeordnungen \(F\) in Beziehung auf \(S\) und \(f\) in Beziehung auf \(s\) wie gleich betrachte, wo denn offenbar auch \(F'\) und \(f'\) gleich sein werden. Es ist hieraus von selbst klar, was in dem am Schluss des zweiten Artikels [S. {\small 408}] erwähnten Falle, wo eine Schicht complexer Grössen durch \textls[120]{mehrere} von einander verschiedene Züge begrenzt wird, unter gleicher Folgeordnung in diesen Zügen zu verstehen ist. Es verhält sich damit ebenso, als wenn die sämmtlichen Uferlinien eines[\textsuperscript{*})] eine oder mehrere Inseln einschliessenden Sees in einem solchen Sinn durchlaufen werden, dass man allemahl den See an derselben Seite, d. i. entweder allemahl rechts oder allemahl links hat.
\begingroup\makeatletter\long\def\@makefntext#1{\noindent\hspace*{1.2em}#1}\makeatother
\footnotetext{[\textsuperscript{*}) Hier stehen in der Handschrift noch die Worte: \textls[120]{Sees, welcher}]}\endgroup
